\documentclass[12pt,letterpaper]{article}
\usepackage[utf8]{inputenc}
\usepackage[spanish]{babel}
\usepackage[margin=2.5cm]{geometry}
\usepackage{amsmath}
\usepackage{amsfonts}
\usepackage{amssymb}

\begin{document}

\begin{center}
    \textbf{\Large Solucionario del Examen de la Unidad II}\\
    \vspace{0.2cm}
    \textbf{Física para Ciencias de la Salud (005-1713)}\\
    \vspace{0.2cm}
    \textit{Cinemática, Dinámica, Trabajo, Energía y Potencia}
\end{center}

\vspace{0.5cm}

\section*{Resolución de los Problemas}

\subsection*{1. Cinemática (Aceleración del flujo sanguíneo)}
\textbf{Datos:} $v_i = 0 \text{ m/s}$ (reposo), $v_f = 0.6 \text{ m/s}$, $t = 0.15 \text{ s}$. \\
\textbf{Fórmula:} $a = \frac{v_f - v_i}{t}$ \\
\textbf{Desarrollo:} 
\[ a = \frac{0.6 \text{ m/s} - 0 \text{ m/s}}{0.15 \text{ s}} = \frac{0.6}{0.15} \text{ m/s}^2 = 4 \text{ m/s}^2 \]
\textbf{Respuesta:} La aceleración media que experimenta la sangre es de \textbf{$4 \text{ m/s}^2$}.

\subsection*{2. Dinámica (Leyes de Newton)}
\textbf{Datos:} $m = 25 \text{ kg}$, $F_{neta} = 150 \text{ N}$. \\
\textbf{Fórmula:} Segunda Ley de Newton, $F = m \cdot a \implies a = \frac{F}{m}$ \\
\textbf{Desarrollo:} 
\[ a = \frac{150 \text{ N}}{25 \text{ kg}} = 6 \text{ m/s}^2 \]
\textbf{Respuesta:} El bloque adquiere una aceleración de \textbf{$6 \text{ m/s}^2$}.

\subsection*{3. Trabajo Mecánico}
\textbf{Datos:} $F = 400 \text{ N}$, $d = 0.8 \text{ m}$. (El ángulo entre la fuerza y el desplazamiento es $0^\circ$ ya que ambos son hacia arriba). \\
\textbf{Fórmula:} $W = F \cdot d \cdot \cos(\theta)$ \\
\textbf{Desarrollo:} 
\[ W = 400 \text{ N} \cdot 0.8 \text{ m} \cdot \cos(0^\circ) = 320 \text{ J} \]
\textbf{Respuesta:} El trabajo mecánico realizado por el paramédico es de \textbf{$320 \text{ Joules}$}.

\subsection*{4. Energía Potencial}
\textbf{Datos:} $m = 1.2 \text{ kg}$, $h = 1.5 \text{ m}$, $g = 9.8 \text{ m/s}^2$. \\
\textbf{Fórmula:} $E_p = m \cdot g \cdot h$ \\
\textbf{Desarrollo:} 
\[ E_p = 1.2 \text{ kg} \cdot 9.8 \text{ m/s}^2 \cdot 1.5 \text{ m} = 17.64 \text{ J} \]
\textbf{Respuesta:} La energía potencial gravitatoria de la bolsa es de \textbf{$17.64 \text{ Joules}$}.

\subsection*{5. Potencia Mecánica}
\textbf{Datos:} $W = 75 \text{ J}$, $t = 60 \text{ s}$. \\
\textbf{Fórmula:} $P = \frac{W}{t}$ \\
\textbf{Desarrollo:} 
\[ P = \frac{75 \text{ J}}{60 \text{ s}} = 1.25 \text{ W} \]
\textbf{Respuesta:} La potencia media desarrollada por el corazón es de \textbf{$1.25 \text{ Watts}$}.

\end{document}